% GENERATED FILE. Edit canonical lessons, then run npm run generate:books.
% canonical-source-hash: fnv1a64:3d9ac44bbbcdffc6
% canonical-lessons: TE-C331-manta, TE-C331-peludu, TE-C331-pranam, TE-C331-raksincu, TE-C331-parugulu-tiyu

\chapter{Fire!}
\label{ch:te-fire}
\hlchaptermodality{331}

\begin{chapteropening}
\textbf{By the end of this chapter:} \emph{I can say a flame, an explosion, a life, to rescue and to run for it.}

Complete the last lesson of chapter 331: I can say a flame, an explosion, a life, to rescue and to run for it.
\end{chapteropening}

\section[\texorpdfstring{\te{మంట} (maṇṭa)}{మంట (maṇṭa)}]{\te{మంట} (maṇṭa) — a flame; a blaze}

\label{lesson:TE-C331-manta}



\begin{quote}
Before the new one: say the Telugu for a news story, then the Telugu for to publish.
\end{quote}

\subsection*{You'll want to know: \te{మంట}}
\textbf{\te{మంట}} — \emph{maṇṭa} — \textquotedblleft{}a flame; a blaze\textquotedblright{}.

\begin{quote}
\te{వంటగదిలో} \te{మంటలు} \te{లేచాయి}.
\end{quote}

\begin{quote}
\emph{vaṇṭagadilō maṇṭalu lēcāyi.}
\end{quote}

\begin{quote}
Flames rose in the kitchen.
\end{quote}

Plural: \te{మంటలు}, \emph{maṇṭalu}.

The plural \te{మంటలు} is the everyday word for a fire that is spreading.

\begin{cousinweb}
\textbf{\te{మంట}} (\emph{maṇṭa}) is native Dravidian: DEDR 4680 lists Telugu \emph{maṇṭa}, \textquotedblleft{}flame, blaze\textquotedblright{}, beside \emph{maṇḍu}, \textquotedblleft{}to burn, blaze\textquotedblright{}, and Tamil \emph{maṇṭu}, \textquotedblleft{}to blaze up\textquotedblright{}.
\end{cousinweb}

\begin{grammarlens}[title={What you've built}]
One more word for telling what happened.
\end{grammarlens}

\subsection*{Your turn}
\begin{itemize}
\raggedright
  \item \emph{Say it:} \emph{maṇṭa}
  \item \emph{Say it:} \emph{maṇṭa}, once more
  \item \emph{Say it:} \emph{vaṇṭagadilō maṇṭalu lēcāyi.} — Flames rose in the kitchen.
  \item \emph{Recall:} say the Telugu for a news story, then the Telugu for to publish, then say \emph{maṇṭa} again
\end{itemize}

\begin{tcolorbox}[breakable,colback=teal!4,colframe=teal!35!black,title={Before you move on}]
What does \te{మంట} mean? (\textquotedblleft{}A flame; a blaze\textquotedblright{}.) Say it once more.
\end{tcolorbox}
\section[\texorpdfstring{\te{పేలుడు} (pēluḍu)}{పేలుడు (pēluḍu)}]{\te{పేలుడు} (pēluḍu) — an explosion; a blast}

\label{lesson:TE-C331-peludu}



\begin{quote}
Before the new one: say the Telugu for to publish, then the Telugu for a flame.
\end{quote}

\subsection*{You'll want to know: \te{పేలుడు}}
\textbf{\te{పేలుడు}} — \emph{pēluḍu} — \textquotedblleft{}an explosion; a blast\textquotedblright{}.

\begin{quote}
\te{పేలుడు} \te{శబ్దం} \te{చాలా} \te{దూరం} \te{వినిపించింది}.
\end{quote}

\begin{quote}
\emph{pēluḍu śabdaṁ cālā dūraṁ vinipiñcindi.}
\end{quote}

\begin{quote}
The sound of the blast was heard far away.
\end{quote}

The noun of \te{పేలు}, to burst: a balloon or a cracker \te{పేలుతుంది}.

\begin{cousinweb}
\textbf{\te{పేలుడు}} (\emph{pēluḍu}) is native Dravidian, the noun of \emph{pēlu}, \textquotedblleft{}to explode, burst, crack\textquotedblright{}, which DEDR 4194 lists with Tamil \emph{piḷ}, \textquotedblleft{}to burst open\textquotedblright{}.
\end{cousinweb}

\begin{grammarlens}[title={What you've built}]
One more word for telling what happened.
\end{grammarlens}

\subsection*{Your turn}
\begin{itemize}
\raggedright
  \item \emph{Say it:} \emph{pēluḍu}
  \item \emph{Say it:} \emph{pēluḍu}, once more
  \item \emph{Say it:} \emph{pēluḍu śabdaṁ cālā dūraṁ vinipiñcindi.} — The sound of the blast was heard far away.
  \item \emph{Recall:} say the Telugu for to publish, then the Telugu for a flame, then say \emph{pēluḍu} again
\end{itemize}

\begin{tcolorbox}[breakable,colback=teal!4,colframe=teal!35!black,title={Before you move on}]
What does \te{పేలుడు} mean? (\textquotedblleft{}An explosion; a blast\textquotedblright{}.) Say it once more.
\end{tcolorbox}
\section[\texorpdfstring{\te{ప్రాణం} (prāṇaṁ)}{ప్రాణం (prāṇaṁ)}]{\te{ప్రాణం} (prāṇaṁ) — a life; the breath of life}

\label{lesson:TE-C331-pranam}



\begin{quote}
Before the new one: say the Telugu for a flame, then the Telugu for an explosion.
\end{quote}

\subsection*{You'll want to know: \te{ప్రాణం}}
\textbf{\te{ప్రాణం}} — \emph{prāṇaṁ} — \textquotedblleft{}a life; the breath of life\textquotedblright{}.

\begin{quote}
\te{ఆ} \te{మంటల్లో} \te{చాలా} \te{ప్రాణాలు} \te{పోయాయి}.
\end{quote}

\begin{quote}
\emph{ā maṇṭallō cālā prāṇālu pōyāyi.}
\end{quote}

\begin{quote}
Many lives were lost in that fire.
\end{quote}

Plural: \te{ప్రాణాలు}, \emph{prāṇālu}.

\te{ప్రాణం} \te{పోయింది}, the life went: a plain way to say that someone died.

\begin{cousinweb}
\textbf{\te{ప్రాణం}} (\emph{prāṇaṁ}) is borrowed from Sanskrit \emph{prāṇa}, \textquotedblleft{}breath, the breath of life\textquotedblright{} (CDIAL 8928).
\end{cousinweb}

\begin{grammarlens}[title={What you've built}]
One more word for telling what happened.
\end{grammarlens}

\subsection*{Your turn}
\begin{itemize}
\raggedright
  \item \emph{Say it:} \emph{prāṇaṁ}
  \item \emph{Say it:} \emph{prāṇaṁ}, once more
  \item \emph{Say it:} \emph{ā maṇṭallō cālā prāṇālu pōyāyi.} — Many lives were lost in that fire.
  \item \emph{Recall:} say the Telugu for a flame, then the Telugu for an explosion, then say \emph{prāṇaṁ} again
\end{itemize}

\begin{tcolorbox}[breakable,colback=teal!4,colframe=teal!35!black,title={Before you move on}]
What does \te{ప్రాణం} mean? (\textquotedblleft{}A life; the breath of life\textquotedblright{}.) Say it once more.
\end{tcolorbox}
\section[\texorpdfstring{\te{రక్షించు} (rakṣiñcu)}{రక్షించు (rakṣiñcu)}]{\te{రక్షించు} (rakṣiñcu) — to rescue; to save}

\label{lesson:TE-C331-raksincu}



\begin{quote}
Before the new one: say the Telugu for an explosion, then the Telugu for a life.
\end{quote}

\subsection*{You'll want to know: \te{రక్షించు}}
\textbf{\te{రక్షించు}} — \emph{rakṣiñcu} — \textquotedblleft{}to rescue; to save\textquotedblright{}.

\begin{quote}
\te{పోలీసులు} \te{పిల్లలను} \te{మంటల} \te{నుంచి} \te{రక్షించారు}.
\end{quote}

\begin{quote}
\emph{pōlīsulu pillalanu maṇṭala nuñci rakṣiñcāru.}
\end{quote}

\begin{quote}
The police rescued the children from the flames.
\end{quote}

Past (I …): \te{రక్షించాను}, \emph{rakṣiñcānu}.

Who is saved takes -\te{ను}; what from takes \te{నుంచి}.

\begin{cousinweb}
\textbf{\te{రక్షించు}} (\emph{rakṣiñcu}) is built on the Sanskrit root \emph{rakṣ}, \textquotedblleft{}to guard, protect\textquotedblright{} (CDIAL 10547), with the Telugu verb ending \emph{-iñcu}.
\end{cousinweb}

\begin{grammarlens}[title={What you've built}]
One more word for telling what happened.
\end{grammarlens}

\subsection*{Your turn}
\begin{itemize}
\raggedright
  \item \emph{Say it:} \emph{rakṣiñcu}
  \item \emph{Say it:} \emph{rakṣiñcu}, once more
  \item \emph{Say it:} \emph{pōlīsulu pillalanu maṇṭala nuñci rakṣiñcāru.} — The police rescued the children from the flames.
  \item \emph{Recall:} say the Telugu for an explosion, then the Telugu for a life, then say \emph{rakṣiñcu} again
\end{itemize}

\begin{tcolorbox}[breakable,colback=teal!4,colframe=teal!35!black,title={Before you move on}]
What does \te{రక్షించు} mean? (\textquotedblleft{}To rescue; to save\textquotedblright{}.) Say it once more.
\end{tcolorbox}
\section[\texorpdfstring{\te{పరుగులు} \te{తీయు} (parugulu tīyu)}{పరుగులు తీయు (parugulu tīyu)}]{\te{పరుగులు} \te{తీయు} (parugulu tīyu) — to run for it; to rush}

\label{lesson:TE-C331-parugulu-tiyu}



\begin{quote}
Before the new one: say the Telugu for a life, then the Telugu for to rescue.
\end{quote}

\subsection*{You'll want to know: \te{పరుగులు} \te{తీయు}}
\textbf{\te{పరుగులు} \te{తీయు}} — \emph{parugulu tīyu} — \textquotedblleft{}to run for it; to rush\textquotedblright{}.

\begin{quote}
\te{మంటలు} \te{చూసి} \te{జనం} \te{పరుగులు} \te{తీశారు}.
\end{quote}

\begin{quote}
\emph{maṇṭalu cūsi janaṁ parugulu tīśāru.}
\end{quote}

\begin{quote}
Seeing the flames, people ran for it.
\end{quote}

Past (I …): \te{పరుగులు} \te{తీశాను}, \emph{parugulu tīśānu}.

Literally to draw runs, with \te{తీయు}, to pull: running in a hurry or in fear.

\begin{cousinweb}
\textbf{\te{పరుగులు} \te{తీయు}} (\emph{parugulu tīyu}) is a native Telugu phrase: \emph{parugu}, \textquotedblleft{}a run\textquotedblright{} (DEDR 3963, with Tamil \emph{pari}, \textquotedblleft{}to run\textquotedblright{}), and \emph{tīyu}, \textquotedblleft{}to draw, pull out\textquotedblright{}.
\end{cousinweb}

\begin{grammarlens}[title={What you've built}]
One more word for telling what happened.
\end{grammarlens}

\subsection*{Your turn}
\begin{itemize}
\raggedright
  \item \emph{Say it:} \emph{parugulu tīyu}
  \item \emph{Say it:} \emph{parugulu tīyu}, once more
  \item \emph{Say it:} \emph{maṇṭalu cūsi janaṁ parugulu tīśāru.} — Seeing the flames, people ran for it.
  \item \emph{Recall:} say the Telugu for a life, then the Telugu for to rescue, then say \emph{parugulu tīyu} again
\end{itemize}

\begin{tcolorbox}[breakable,colback=teal!4,colframe=teal!35!black,title={Before you move on}]
What does \te{పరుగులు} \te{తీయు} mean? (\textquotedblleft{}To run for it; to rush\textquotedblright{}.) Say it once more.
\end{tcolorbox}
